\documentclass[11pt]{article}

    \usepackage[breakable]{tcolorbox}
    \usepackage{parskip} % Stop auto-indenting (to mimic markdown behaviour)
    

    % Basic figure setup, for now with no caption control since it's done
    % automatically by Pandoc (which extracts ![](path) syntax from Markdown).
    \usepackage{graphicx}
    % Keep aspect ratio if custom image width or height is specified
    \setkeys{Gin}{keepaspectratio}
    % Maintain compatibility with old templates. Remove in nbconvert 6.0
    \let\Oldincludegraphics\includegraphics
    % Ensure that by default, figures have no caption (until we provide a
    % proper Figure object with a Caption API and a way to capture that
    % in the conversion process - todo).
    \usepackage{caption}
    \DeclareCaptionFormat{nocaption}{}
    \captionsetup{format=nocaption,aboveskip=0pt,belowskip=0pt}

    \usepackage{float}
    \floatplacement{figure}{H} % forces figures to be placed at the correct location
    \usepackage{xcolor} % Allow colors to be defined
    \usepackage{enumerate} % Needed for markdown enumerations to work
    \usepackage{geometry} % Used to adjust the document margins
    \usepackage{amsmath} % Equations
    \usepackage{amssymb} % Equations
    \usepackage{textcomp} % defines textquotesingle
    % Hack from http://tex.stackexchange.com/a/47451/13684:
    \AtBeginDocument{%
        \def\PYZsq{\textquotesingle}% Upright quotes in Pygmentized code
    }
    \usepackage{upquote} % Upright quotes for verbatim code
    \usepackage{eurosym} % defines \euro

    \usepackage{iftex}
    \ifPDFTeX
        \usepackage[T1]{fontenc}
        \IfFileExists{alphabeta.sty}{
              \usepackage{alphabeta}
          }{
              \usepackage[mathletters]{ucs}
              \usepackage[utf8x]{inputenc}
          }
    \else
        \usepackage{fontspec}
        \usepackage{unicode-math}
    \fi

    \usepackage{fancyvrb} % verbatim replacement that allows latex
    \usepackage{grffile} % extends the file name processing of package graphics
                         % to support a larger range
    \makeatletter % fix for old versions of grffile with XeLaTeX
    \@ifpackagelater{grffile}{2019/11/01}
    {
      % Do nothing on new versions
    }
    {
      \def\Gread@@xetex#1{%
        \IfFileExists{"\Gin@base".bb}%
        {\Gread@eps{\Gin@base.bb}}%
        {\Gread@@xetex@aux#1}%
      }
    }
    \makeatother
    \usepackage[Export]{adjustbox} % Used to constrain images to a maximum size
    \adjustboxset{max size={0.9\linewidth}{0.9\paperheight}}

    % The hyperref package gives us a pdf with properly built
    % internal navigation ('pdf bookmarks' for the table of contents,
    % internal cross-reference links, web links for URLs, etc.)
    \usepackage{hyperref}
    % The default LaTeX title has an obnoxious amount of whitespace. By default,
    % titling removes some of it. It also provides customization options.
    \usepackage{titling}
    \usepackage{longtable} % longtable support required by pandoc >1.10
    \usepackage{booktabs}  % table support for pandoc > 1.12.2
    \usepackage{array}     % table support for pandoc >= 2.11.3
    \usepackage{calc}      % table minipage width calculation for pandoc >= 2.11.1
    \usepackage[inline]{enumitem} % IRkernel/repr support (it uses the enumerate* environment)
    \usepackage[normalem]{ulem} % ulem is needed to support strikethroughs (\sout)
                                % normalem makes italics be italics, not underlines
    \usepackage{soul}      % strikethrough (\st) support for pandoc >= 3.0.0
    \usepackage{mathrsfs}
    

    
    % Colors for the hyperref package
    \definecolor{urlcolor}{rgb}{0,.145,.698}
    \definecolor{linkcolor}{rgb}{.71,0.21,0.01}
    \definecolor{citecolor}{rgb}{.12,.54,.11}

    % ANSI colors
    \definecolor{ansi-black}{HTML}{3E424D}
    \definecolor{ansi-black-intense}{HTML}{282C36}
    \definecolor{ansi-red}{HTML}{E75C58}
    \definecolor{ansi-red-intense}{HTML}{B22B31}
    \definecolor{ansi-green}{HTML}{00A250}
    \definecolor{ansi-green-intense}{HTML}{007427}
    \definecolor{ansi-yellow}{HTML}{DDB62B}
    \definecolor{ansi-yellow-intense}{HTML}{B27D12}
    \definecolor{ansi-blue}{HTML}{208FFB}
    \definecolor{ansi-blue-intense}{HTML}{0065CA}
    \definecolor{ansi-magenta}{HTML}{D160C4}
    \definecolor{ansi-magenta-intense}{HTML}{A03196}
    \definecolor{ansi-cyan}{HTML}{60C6C8}
    \definecolor{ansi-cyan-intense}{HTML}{258F8F}
    \definecolor{ansi-white}{HTML}{C5C1B4}
    \definecolor{ansi-white-intense}{HTML}{A1A6B2}
    \definecolor{ansi-default-inverse-fg}{HTML}{FFFFFF}
    \definecolor{ansi-default-inverse-bg}{HTML}{000000}

    % common color for the border for error outputs.
    \definecolor{outerrorbackground}{HTML}{FFDFDF}

    % commands and environments needed by pandoc snippets
    % extracted from the output of `pandoc -s`
    \providecommand{\tightlist}{%
      \setlength{\itemsep}{0pt}\setlength{\parskip}{0pt}}
    \DefineVerbatimEnvironment{Highlighting}{Verbatim}{commandchars=\\\{\}}
    % Add ',fontsize=\small' for more characters per line
    \newenvironment{Shaded}{}{}
    \newcommand{\KeywordTok}[1]{\textcolor[rgb]{0.00,0.44,0.13}{\textbf{{#1}}}}
    \newcommand{\DataTypeTok}[1]{\textcolor[rgb]{0.56,0.13,0.00}{{#1}}}
    \newcommand{\DecValTok}[1]{\textcolor[rgb]{0.25,0.63,0.44}{{#1}}}
    \newcommand{\BaseNTok}[1]{\textcolor[rgb]{0.25,0.63,0.44}{{#1}}}
    \newcommand{\FloatTok}[1]{\textcolor[rgb]{0.25,0.63,0.44}{{#1}}}
    \newcommand{\CharTok}[1]{\textcolor[rgb]{0.25,0.44,0.63}{{#1}}}
    \newcommand{\StringTok}[1]{\textcolor[rgb]{0.25,0.44,0.63}{{#1}}}
    \newcommand{\CommentTok}[1]{\textcolor[rgb]{0.38,0.63,0.69}{\textit{{#1}}}}
    \newcommand{\OtherTok}[1]{\textcolor[rgb]{0.00,0.44,0.13}{{#1}}}
    \newcommand{\AlertTok}[1]{\textcolor[rgb]{1.00,0.00,0.00}{\textbf{{#1}}}}
    \newcommand{\FunctionTok}[1]{\textcolor[rgb]{0.02,0.16,0.49}{{#1}}}
    \newcommand{\RegionMarkerTok}[1]{{#1}}
    \newcommand{\ErrorTok}[1]{\textcolor[rgb]{1.00,0.00,0.00}{\textbf{{#1}}}}
    \newcommand{\NormalTok}[1]{{#1}}

    % Additional commands for more recent versions of Pandoc
    \newcommand{\ConstantTok}[1]{\textcolor[rgb]{0.53,0.00,0.00}{{#1}}}
    \newcommand{\SpecialCharTok}[1]{\textcolor[rgb]{0.25,0.44,0.63}{{#1}}}
    \newcommand{\VerbatimStringTok}[1]{\textcolor[rgb]{0.25,0.44,0.63}{{#1}}}
    \newcommand{\SpecialStringTok}[1]{\textcolor[rgb]{0.73,0.40,0.53}{{#1}}}
    \newcommand{\ImportTok}[1]{{#1}}
    \newcommand{\DocumentationTok}[1]{\textcolor[rgb]{0.73,0.13,0.13}{\textit{{#1}}}}
    \newcommand{\AnnotationTok}[1]{\textcolor[rgb]{0.38,0.63,0.69}{\textbf{\textit{{#1}}}}}
    \newcommand{\CommentVarTok}[1]{\textcolor[rgb]{0.38,0.63,0.69}{\textbf{\textit{{#1}}}}}
    \newcommand{\VariableTok}[1]{\textcolor[rgb]{0.10,0.09,0.49}{{#1}}}
    \newcommand{\ControlFlowTok}[1]{\textcolor[rgb]{0.00,0.44,0.13}{\textbf{{#1}}}}
    \newcommand{\OperatorTok}[1]{\textcolor[rgb]{0.40,0.40,0.40}{{#1}}}
    \newcommand{\BuiltInTok}[1]{{#1}}
    \newcommand{\ExtensionTok}[1]{{#1}}
    \newcommand{\PreprocessorTok}[1]{\textcolor[rgb]{0.74,0.48,0.00}{{#1}}}
    \newcommand{\AttributeTok}[1]{\textcolor[rgb]{0.49,0.56,0.16}{{#1}}}
    \newcommand{\InformationTok}[1]{\textcolor[rgb]{0.38,0.63,0.69}{\textbf{\textit{{#1}}}}}
    \newcommand{\WarningTok}[1]{\textcolor[rgb]{0.38,0.63,0.69}{\textbf{\textit{{#1}}}}}
    \makeatletter
    \newsavebox\pandoc@box
    \newcommand*\pandocbounded[1]{%
      \sbox\pandoc@box{#1}%
      % scaling factors for width and height
      \Gscale@div\@tempa\textheight{\dimexpr\ht\pandoc@box+\dp\pandoc@box\relax}%
      \Gscale@div\@tempb\linewidth{\wd\pandoc@box}%
      % select the smaller of both
      \ifdim\@tempb\p@<\@tempa\p@
        \let\@tempa\@tempb
      \fi
      % scaling accordingly (\@tempa < 1)
      \ifdim\@tempa\p@<\p@
        \scalebox{\@tempa}{\usebox\pandoc@box}%
      % scaling not needed, use as it is
      \else
        \usebox{\pandoc@box}%
      \fi
    }
    \makeatother

    % Define a nice break command that doesn't care if a line doesn't already
    % exist.
    \def\br{\hspace*{\fill} \\* }
    % Math Jax compatibility definitions
    \def\gt{>}
    \def\lt{<}
    \let\Oldtex\TeX
    \let\Oldlatex\LaTeX
    \renewcommand{\TeX}{\textrm{\Oldtex}}
    \renewcommand{\LaTeX}{\textrm{\Oldlatex}}
    % Document parameters
    % Document title
    \title{587-N15b}
    
    
    
    
    
    
    
% Pygments definitions
\makeatletter
\def\PY@reset{\let\PY@it=\relax \let\PY@bf=\relax%
    \let\PY@ul=\relax \let\PY@tc=\relax%
    \let\PY@bc=\relax \let\PY@ff=\relax}
\def\PY@tok#1{\csname PY@tok@#1\endcsname}
\def\PY@toks#1+{\ifx\relax#1\empty\else%
    \PY@tok{#1}\expandafter\PY@toks\fi}
\def\PY@do#1{\PY@bc{\PY@tc{\PY@ul{%
    \PY@it{\PY@bf{\PY@ff{#1}}}}}}}
\def\PY#1#2{\PY@reset\PY@toks#1+\relax+\PY@do{#2}}

\@namedef{PY@tok@w}{\def\PY@tc##1{\textcolor[rgb]{0.73,0.73,0.73}{##1}}}
\@namedef{PY@tok@c}{\let\PY@it=\textit\def\PY@tc##1{\textcolor[rgb]{0.24,0.48,0.48}{##1}}}
\@namedef{PY@tok@cp}{\def\PY@tc##1{\textcolor[rgb]{0.61,0.40,0.00}{##1}}}
\@namedef{PY@tok@k}{\let\PY@bf=\textbf\def\PY@tc##1{\textcolor[rgb]{0.00,0.50,0.00}{##1}}}
\@namedef{PY@tok@kp}{\def\PY@tc##1{\textcolor[rgb]{0.00,0.50,0.00}{##1}}}
\@namedef{PY@tok@kt}{\def\PY@tc##1{\textcolor[rgb]{0.69,0.00,0.25}{##1}}}
\@namedef{PY@tok@o}{\def\PY@tc##1{\textcolor[rgb]{0.40,0.40,0.40}{##1}}}
\@namedef{PY@tok@ow}{\let\PY@bf=\textbf\def\PY@tc##1{\textcolor[rgb]{0.67,0.13,1.00}{##1}}}
\@namedef{PY@tok@nb}{\def\PY@tc##1{\textcolor[rgb]{0.00,0.50,0.00}{##1}}}
\@namedef{PY@tok@nf}{\def\PY@tc##1{\textcolor[rgb]{0.00,0.00,1.00}{##1}}}
\@namedef{PY@tok@nc}{\let\PY@bf=\textbf\def\PY@tc##1{\textcolor[rgb]{0.00,0.00,1.00}{##1}}}
\@namedef{PY@tok@nn}{\let\PY@bf=\textbf\def\PY@tc##1{\textcolor[rgb]{0.00,0.00,1.00}{##1}}}
\@namedef{PY@tok@ne}{\let\PY@bf=\textbf\def\PY@tc##1{\textcolor[rgb]{0.80,0.25,0.22}{##1}}}
\@namedef{PY@tok@nv}{\def\PY@tc##1{\textcolor[rgb]{0.10,0.09,0.49}{##1}}}
\@namedef{PY@tok@no}{\def\PY@tc##1{\textcolor[rgb]{0.53,0.00,0.00}{##1}}}
\@namedef{PY@tok@nl}{\def\PY@tc##1{\textcolor[rgb]{0.46,0.46,0.00}{##1}}}
\@namedef{PY@tok@ni}{\let\PY@bf=\textbf\def\PY@tc##1{\textcolor[rgb]{0.44,0.44,0.44}{##1}}}
\@namedef{PY@tok@na}{\def\PY@tc##1{\textcolor[rgb]{0.41,0.47,0.13}{##1}}}
\@namedef{PY@tok@nt}{\let\PY@bf=\textbf\def\PY@tc##1{\textcolor[rgb]{0.00,0.50,0.00}{##1}}}
\@namedef{PY@tok@nd}{\def\PY@tc##1{\textcolor[rgb]{0.67,0.13,1.00}{##1}}}
\@namedef{PY@tok@s}{\def\PY@tc##1{\textcolor[rgb]{0.73,0.13,0.13}{##1}}}
\@namedef{PY@tok@sd}{\let\PY@it=\textit\def\PY@tc##1{\textcolor[rgb]{0.73,0.13,0.13}{##1}}}
\@namedef{PY@tok@si}{\let\PY@bf=\textbf\def\PY@tc##1{\textcolor[rgb]{0.64,0.35,0.47}{##1}}}
\@namedef{PY@tok@se}{\let\PY@bf=\textbf\def\PY@tc##1{\textcolor[rgb]{0.67,0.36,0.12}{##1}}}
\@namedef{PY@tok@sr}{\def\PY@tc##1{\textcolor[rgb]{0.64,0.35,0.47}{##1}}}
\@namedef{PY@tok@ss}{\def\PY@tc##1{\textcolor[rgb]{0.10,0.09,0.49}{##1}}}
\@namedef{PY@tok@sx}{\def\PY@tc##1{\textcolor[rgb]{0.00,0.50,0.00}{##1}}}
\@namedef{PY@tok@m}{\def\PY@tc##1{\textcolor[rgb]{0.40,0.40,0.40}{##1}}}
\@namedef{PY@tok@gh}{\let\PY@bf=\textbf\def\PY@tc##1{\textcolor[rgb]{0.00,0.00,0.50}{##1}}}
\@namedef{PY@tok@gu}{\let\PY@bf=\textbf\def\PY@tc##1{\textcolor[rgb]{0.50,0.00,0.50}{##1}}}
\@namedef{PY@tok@gd}{\def\PY@tc##1{\textcolor[rgb]{0.63,0.00,0.00}{##1}}}
\@namedef{PY@tok@gi}{\def\PY@tc##1{\textcolor[rgb]{0.00,0.52,0.00}{##1}}}
\@namedef{PY@tok@gr}{\def\PY@tc##1{\textcolor[rgb]{0.89,0.00,0.00}{##1}}}
\@namedef{PY@tok@ge}{\let\PY@it=\textit}
\@namedef{PY@tok@gs}{\let\PY@bf=\textbf}
\@namedef{PY@tok@ges}{\let\PY@bf=\textbf\let\PY@it=\textit}
\@namedef{PY@tok@gp}{\let\PY@bf=\textbf\def\PY@tc##1{\textcolor[rgb]{0.00,0.00,0.50}{##1}}}
\@namedef{PY@tok@go}{\def\PY@tc##1{\textcolor[rgb]{0.44,0.44,0.44}{##1}}}
\@namedef{PY@tok@gt}{\def\PY@tc##1{\textcolor[rgb]{0.00,0.27,0.87}{##1}}}
\@namedef{PY@tok@err}{\def\PY@bc##1{{\setlength{\fboxsep}{\string -\fboxrule}\fcolorbox[rgb]{1.00,0.00,0.00}{1,1,1}{\strut ##1}}}}
\@namedef{PY@tok@kc}{\let\PY@bf=\textbf\def\PY@tc##1{\textcolor[rgb]{0.00,0.50,0.00}{##1}}}
\@namedef{PY@tok@kd}{\let\PY@bf=\textbf\def\PY@tc##1{\textcolor[rgb]{0.00,0.50,0.00}{##1}}}
\@namedef{PY@tok@kn}{\let\PY@bf=\textbf\def\PY@tc##1{\textcolor[rgb]{0.00,0.50,0.00}{##1}}}
\@namedef{PY@tok@kr}{\let\PY@bf=\textbf\def\PY@tc##1{\textcolor[rgb]{0.00,0.50,0.00}{##1}}}
\@namedef{PY@tok@bp}{\def\PY@tc##1{\textcolor[rgb]{0.00,0.50,0.00}{##1}}}
\@namedef{PY@tok@fm}{\def\PY@tc##1{\textcolor[rgb]{0.00,0.00,1.00}{##1}}}
\@namedef{PY@tok@vc}{\def\PY@tc##1{\textcolor[rgb]{0.10,0.09,0.49}{##1}}}
\@namedef{PY@tok@vg}{\def\PY@tc##1{\textcolor[rgb]{0.10,0.09,0.49}{##1}}}
\@namedef{PY@tok@vi}{\def\PY@tc##1{\textcolor[rgb]{0.10,0.09,0.49}{##1}}}
\@namedef{PY@tok@vm}{\def\PY@tc##1{\textcolor[rgb]{0.10,0.09,0.49}{##1}}}
\@namedef{PY@tok@sa}{\def\PY@tc##1{\textcolor[rgb]{0.73,0.13,0.13}{##1}}}
\@namedef{PY@tok@sb}{\def\PY@tc##1{\textcolor[rgb]{0.73,0.13,0.13}{##1}}}
\@namedef{PY@tok@sc}{\def\PY@tc##1{\textcolor[rgb]{0.73,0.13,0.13}{##1}}}
\@namedef{PY@tok@dl}{\def\PY@tc##1{\textcolor[rgb]{0.73,0.13,0.13}{##1}}}
\@namedef{PY@tok@s2}{\def\PY@tc##1{\textcolor[rgb]{0.73,0.13,0.13}{##1}}}
\@namedef{PY@tok@sh}{\def\PY@tc##1{\textcolor[rgb]{0.73,0.13,0.13}{##1}}}
\@namedef{PY@tok@s1}{\def\PY@tc##1{\textcolor[rgb]{0.73,0.13,0.13}{##1}}}
\@namedef{PY@tok@mb}{\def\PY@tc##1{\textcolor[rgb]{0.40,0.40,0.40}{##1}}}
\@namedef{PY@tok@mf}{\def\PY@tc##1{\textcolor[rgb]{0.40,0.40,0.40}{##1}}}
\@namedef{PY@tok@mh}{\def\PY@tc##1{\textcolor[rgb]{0.40,0.40,0.40}{##1}}}
\@namedef{PY@tok@mi}{\def\PY@tc##1{\textcolor[rgb]{0.40,0.40,0.40}{##1}}}
\@namedef{PY@tok@il}{\def\PY@tc##1{\textcolor[rgb]{0.40,0.40,0.40}{##1}}}
\@namedef{PY@tok@mo}{\def\PY@tc##1{\textcolor[rgb]{0.40,0.40,0.40}{##1}}}
\@namedef{PY@tok@ch}{\let\PY@it=\textit\def\PY@tc##1{\textcolor[rgb]{0.24,0.48,0.48}{##1}}}
\@namedef{PY@tok@cm}{\let\PY@it=\textit\def\PY@tc##1{\textcolor[rgb]{0.24,0.48,0.48}{##1}}}
\@namedef{PY@tok@cpf}{\let\PY@it=\textit\def\PY@tc##1{\textcolor[rgb]{0.24,0.48,0.48}{##1}}}
\@namedef{PY@tok@c1}{\let\PY@it=\textit\def\PY@tc##1{\textcolor[rgb]{0.24,0.48,0.48}{##1}}}
\@namedef{PY@tok@cs}{\let\PY@it=\textit\def\PY@tc##1{\textcolor[rgb]{0.24,0.48,0.48}{##1}}}

\def\PYZbs{\char`\\}
\def\PYZus{\char`\_}
\def\PYZob{\char`\{}
\def\PYZcb{\char`\}}
\def\PYZca{\char`\^}
\def\PYZam{\char`\&}
\def\PYZlt{\char`\<}
\def\PYZgt{\char`\>}
\def\PYZsh{\char`\#}
\def\PYZpc{\char`\%}
\def\PYZdl{\char`\$}
\def\PYZhy{\char`\-}
\def\PYZsq{\char`\'}
\def\PYZdq{\char`\"}
\def\PYZti{\char`\~}
% for compatibility with earlier versions
\def\PYZat{@}
\def\PYZlb{[}
\def\PYZrb{]}
\makeatother


    % For linebreaks inside Verbatim environment from package fancyvrb.
    \makeatletter
        \newbox\Wrappedcontinuationbox
        \newbox\Wrappedvisiblespacebox
        \newcommand*\Wrappedvisiblespace {\textcolor{red}{\textvisiblespace}}
        \newcommand*\Wrappedcontinuationsymbol {\textcolor{red}{\llap{\tiny$\m@th\hookrightarrow$}}}
        \newcommand*\Wrappedcontinuationindent {3ex }
        \newcommand*\Wrappedafterbreak {\kern\Wrappedcontinuationindent\copy\Wrappedcontinuationbox}
        % Take advantage of the already applied Pygments mark-up to insert
        % potential linebreaks for TeX processing.
        %        {, <, #, %, $, ' and ": go to next line.
        %        _, }, ^, &, >, - and ~: stay at end of broken line.
        % Use of \textquotesingle for straight quote.
        \newcommand*\Wrappedbreaksatspecials {%
            \def\PYGZus{\discretionary{\char`\_}{\Wrappedafterbreak}{\char`\_}}%
            \def\PYGZob{\discretionary{}{\Wrappedafterbreak\char`\{}{\char`\{}}%
            \def\PYGZcb{\discretionary{\char`\}}{\Wrappedafterbreak}{\char`\}}}%
            \def\PYGZca{\discretionary{\char`\^}{\Wrappedafterbreak}{\char`\^}}%
            \def\PYGZam{\discretionary{\char`\&}{\Wrappedafterbreak}{\char`\&}}%
            \def\PYGZlt{\discretionary{}{\Wrappedafterbreak\char`\<}{\char`\<}}%
            \def\PYGZgt{\discretionary{\char`\>}{\Wrappedafterbreak}{\char`\>}}%
            \def\PYGZsh{\discretionary{}{\Wrappedafterbreak\char`\#}{\char`\#}}%
            \def\PYGZpc{\discretionary{}{\Wrappedafterbreak\char`\%}{\char`\%}}%
            \def\PYGZdl{\discretionary{}{\Wrappedafterbreak\char`\$}{\char`\$}}%
            \def\PYGZhy{\discretionary{\char`\-}{\Wrappedafterbreak}{\char`\-}}%
            \def\PYGZsq{\discretionary{}{\Wrappedafterbreak\textquotesingle}{\textquotesingle}}%
            \def\PYGZdq{\discretionary{}{\Wrappedafterbreak\char`\"}{\char`\"}}%
            \def\PYGZti{\discretionary{\char`\~}{\Wrappedafterbreak}{\char`\~}}%
        }
        % Some characters . , ; ? ! / are not pygmentized.
        % This macro makes them "active" and they will insert potential linebreaks
        \newcommand*\Wrappedbreaksatpunct {%
            \lccode`\~`\.\lowercase{\def~}{\discretionary{\hbox{\char`\.}}{\Wrappedafterbreak}{\hbox{\char`\.}}}%
            \lccode`\~`\,\lowercase{\def~}{\discretionary{\hbox{\char`\,}}{\Wrappedafterbreak}{\hbox{\char`\,}}}%
            \lccode`\~`\;\lowercase{\def~}{\discretionary{\hbox{\char`\;}}{\Wrappedafterbreak}{\hbox{\char`\;}}}%
            \lccode`\~`\:\lowercase{\def~}{\discretionary{\hbox{\char`\:}}{\Wrappedafterbreak}{\hbox{\char`\:}}}%
            \lccode`\~`\?\lowercase{\def~}{\discretionary{\hbox{\char`\?}}{\Wrappedafterbreak}{\hbox{\char`\?}}}%
            \lccode`\~`\!\lowercase{\def~}{\discretionary{\hbox{\char`\!}}{\Wrappedafterbreak}{\hbox{\char`\!}}}%
            \lccode`\~`\/\lowercase{\def~}{\discretionary{\hbox{\char`\/}}{\Wrappedafterbreak}{\hbox{\char`\/}}}%
            \catcode`\.\active
            \catcode`\,\active
            \catcode`\;\active
            \catcode`\:\active
            \catcode`\?\active
            \catcode`\!\active
            \catcode`\/\active
            \lccode`\~`\~
        }
    \makeatother

    \let\OriginalVerbatim=\Verbatim
    \makeatletter
    \renewcommand{\Verbatim}[1][1]{%
        %\parskip\z@skip
        \sbox\Wrappedcontinuationbox {\Wrappedcontinuationsymbol}%
        \sbox\Wrappedvisiblespacebox {\FV@SetupFont\Wrappedvisiblespace}%
        \def\FancyVerbFormatLine ##1{\hsize\linewidth
            \vtop{\raggedright\hyphenpenalty\z@\exhyphenpenalty\z@
                \doublehyphendemerits\z@\finalhyphendemerits\z@
                \strut ##1\strut}%
        }%
        % If the linebreak is at a space, the latter will be displayed as visible
        % space at end of first line, and a continuation symbol starts next line.
        % Stretch/shrink are however usually zero for typewriter font.
        \def\FV@Space {%
            \nobreak\hskip\z@ plus\fontdimen3\font minus\fontdimen4\font
            \discretionary{\copy\Wrappedvisiblespacebox}{\Wrappedafterbreak}
            {\kern\fontdimen2\font}%
        }%

        % Allow breaks at special characters using \PYG... macros.
        \Wrappedbreaksatspecials
        % Breaks at punctuation characters . , ; ? ! and / need catcode=\active
        \OriginalVerbatim[#1,codes*=\Wrappedbreaksatpunct]%
    }
    \makeatother

    % Exact colors from NB
    \definecolor{incolor}{HTML}{303F9F}
    \definecolor{outcolor}{HTML}{D84315}
    \definecolor{cellborder}{HTML}{CFCFCF}
    \definecolor{cellbackground}{HTML}{F7F7F7}

    % prompt
    \makeatletter
    \newcommand{\boxspacing}{\kern\kvtcb@left@rule\kern\kvtcb@boxsep}
    \makeatother
    \newcommand{\prompt}[4]{
        {\ttfamily\llap{{\color{#2}[#3]:\hspace{3pt}#4}}\vspace{-\baselineskip}}
    }
    

    
    % Prevent overflowing lines due to hard-to-break entities
    \sloppy
    % Setup hyperref package
    \hypersetup{
      breaklinks=true,  % so long urls are correctly broken across lines
      colorlinks=true,
      urlcolor=urlcolor,
      linkcolor=linkcolor,
      citecolor=citecolor,
      }
    % Slightly bigger margins than the latex defaults
    
    \geometry{verbose,tmargin=1in,bmargin=1in,lmargin=1in,rmargin=1in}
    
    

\begin{document}
    
    \maketitle
    
    

    
    \hypertarget{material-balances-iii}{%
\section{Material Balances III}\label{material-balances-iii}}

\href{https://www.youtube.com/watch?v=TSM0jj-hZtw\&list=RDTSM0jj-hZtw\&start_radio=1}{Bond's
favorite Zep tune.}

This lecture applies the general species balance to three common reactor
models: The Batch Reactor, the Continuous Stirred Tank Reactor (CSTR),
and the Plug Flow Reactor (PFR).

    \hypertarget{the-general-balance-equation}{%
\subsection{The General Balance
Equation}\label{the-general-balance-equation}}

Regardless of the reactor we're using, we can always apply the general
balance equation that we developed in
\href{https://jqbond.github.io/CEN587/Notebooks/587-N13.html\#developing-a-material-balance-for-a-reacting-system}{Notebook
13}:

\[\frac{dN_j}{dt} = F_{j,f} - F_j + G_j\]

For homogeneous reactions (reactions occuring in a single phase), we
usually report \emph{intensive reaction rates} per unit volume, which
leads to the following result for homogeneous reactions (Developed in
\href{https://jqbond.github.io/CEN587/Notebooks/587-N14.html\#doing-this-for-a-real-reactor}{Notebook
14}):

\[\frac{dN_j}{dt} = F_{j,f} - F_j + \int^V R_jV\]

We further recall that \(R_j\) is the net, intensive production rate of
species \(j\) given in units of moles per volume per time. We calculate
the net intensive production rate by
\href{https://jqbond.github.io/CEN587/Notebooks/587-N14.html\#multiple-reactions}{summing
the production rate of \(j\) in each individual reaction occuring in the
system}:

\[R_j = \sum_{i = 1}^{N_R} \nu_{i,j} \, r_i\]

    \hypertarget{the-batch-reactor}{%
\subsection{The Batch Reactor}\label{the-batch-reactor}}

A batch reactor is a closed system. It is initially charged with some
quantity of reactants, products, solvents, catalysts, etc. at some
temperature, pressure, and composition. Thereafter, there is no flow of
material into or out of the reactor (although exchange of energy with
the surroundings is permitted). Typically, we imagine starting a
reaction at \(t = 0\), and then we allow the reaction to occur as time
progresses. This is not a flow process, and batch reactors do not
operate at steady state (steady state in a batch reactor means that the
reaction has reached completion or equilibrium such that composition is
no longer changing). When we discuss batch reactors, we usually will
think about how the characteristics of the batch reactor and the
reaction media inside of the batch reactor \textbf{\emph{change as a
function of time}}. Mass conservation requires that, for a closed
system, the total mass inside of the system is invariant with time.
Element conservation also applies, so the atomic quantities of elements
will also not change with time. However, as the reaction occurs, the
nature of that mass and elements will change since reactants are
generally converting into products. In fact, for a batch reactor, total
mass and atomic quantities of elements are pretty much the only things
that are fixed. When working with a batch reactor, as time elapses, we
should generally expect that:

\begin{enumerate}
\def\labelenumi{\arabic{enumi}.}
\tightlist
\item
  The number of moles of each species may change with time.
\item
  The total moles inside the reactor may change with time.
\item
  The composition inside the reactor may change with time.
\item
  The temperature inside the reactor may change with time.
\item
  The pressure of the reactor may change with time.
\item
  The size (volume) of the reactor may change with time.
\end{enumerate}

When we analyze batch reactors, we'll always be considering how their
characteristics change with time. During the useful part of their
operation, batch reactions are always dynamic, non-steady state
processes.

\hypertarget{applications-of-batch-reactors}{%
\subsubsection{Applications of Batch
Reactors}\label{applications-of-batch-reactors}}

Batch reactors tend to be a bit more labor intensive since they require
startup, shutdown, product recovery, cleaning/decontamination,
re-charging, etc. They are commonly used for exploratory research with
liquid- or gas-phase reactions at small scales. Often, we study kinetics
in the laboratory using batch reactors. This allows us to develop
functions that tell us how intensive rates of reactions vary with
temperature, pressure, and composition. Once we have developed these
\textbf{\emph{rate laws}}, they are applicable universally, and we can
use them to scale the process to larger reactors. Batch reactors are
also used industrially for the production of things that do not scale
well to continuous flow. Some examples might be very high value products
(like pharmaceuticals) that are produced in very small quantities. In
these cases, it may not be practical to run a process continuously, and
the industry is better served by using a batch reactor to produce small
quantities in multiple batches throughout the year. Processes that use
expensive homogeneous catalysts (organometallic complexes, enzymes,
etc.) are also sometimes carried out in batch processes because it is
impractical to continuously flow expensive catalysts in solution and
subsequently try to recover them from the reactor effluent. In these
cases, it is may be better to work in batch systems. Another example
could be biological processes, such as those involving microbial
fermentation. Fermentation in the food and beverage industry,
beer-brewing, for example is usually performed in batch reactors.

    \hypertarget{material-balance-on-the-batch-reactor}{%
\subsubsection{Material Balance on The Batch
Reactor}\label{material-balance-on-the-batch-reactor}}

A typical batch reactor is illustrated in the figure below. In this
case, we have a homogeneous reaction occuring throughout the fluid phase
(blue color)

Now we'll apply the balance equation to that reactor.

\[\frac{dN_j}{dt} = F_{j,f} - F_j + G_j\]

A batch reactor is a closed system, so it is clear that we do not have
flow of species \(j\) into or out of the reactor. Accordingly, both of
the flow terms in our general balance equation are zero.

\[F_{j,f} = F_j = 0\]

We generally expect that we have a reaction occuring in this reactor, so
that means that there is probably generation or consumption of species
\(j\) by reaction. This means that the generation term is usually
non-zero:

\[G_j \neq 0\]

Finally, we should usually expect that the quantity of each species is
changing inside of a batch reactor as time elapses. Even though the
total mass in the reactor is constant, we should be converting reactants
into products, so there should be a change in the number of moles of
each species with time (otherwise, we have not designed the reactor very
well!!) As discussed in
\href{https://jqbond.github.io/CEN587/Lectures/587-L10.html\#accumulation-without-flow}{Notebook
13}, this is a type of accumulation. Specifically, the fact that the
number of moles of species is generally going to change as a function of
time in a batch reactor means that there is a non-zero accumulation term
in a batch reactor material balance:

\[\frac{dN_j}{dt} \neq 0\]

Based on these insights, we have the following equation for balance on
species \(j\) in a generic batch reactor:

\[\frac{dN_j}{dt} = G_j\]

In other words, the rate of accumulation of species in a batch reactor
is exactly equal to its rate of generation or consumption by chemical
reaction! We can go further. Since this example specifically considers a
homogeneous reaction occuring throughout the volume of fluid occupied by
the reactor, intensive rates will be specified per unit volume, so:

\[\frac{dN_j}{dt} = \int^V R_jdV\]

This equation applies for any batch reactor without making assumptions.
If we want to simplify further, we have to start making assumptions
about the operation of our reactor.

\hypertarget{simplifiying-the-material-balance-for-a-well-mixed-batch-reactor}{%
\subsubsection{Simplifiying the material balance for a well-mixed batch
reactor}\label{simplifiying-the-material-balance-for-a-well-mixed-batch-reactor}}

To facilitate solution of problems in this course, we will always assume
that our batch reactors have \textbf{\emph{perfect mixing}}. This
implies that temperature, pressure, and composition \textbf{\emph{must
be spatially uniform}} inside of the batch reactor. If we could examine
a fluid sample from one position inside of a \textbf{\emph{perfectly
mixed batch reactor}}, we would find that the temperature, pressure, and
composition is exactly the same as the temperature, pressure, and
composition within fluid samples taken \textbf{\emph{at every other
position inside of the reactor}}. This allows us to make a very useful
simplification to the batch reactor material balance. Specifically, we
know that reaction rates are entirely determined by temperature,
pressure, and composition. If those three things are spatially uniform
throughout the batch reactor (i.e., they do not change as a function of
position inside of the reactor), then reaction rates must also be
spatially uniform throughout the batch reactor. If reaction rates are
position independent (spatially invariant), then intensive species
production rates are also position independent:

\[R_j \neq f(V)\]

If we're able to assume that \(R_j\) is not a function of volume, then
the volume integral in our material balance simplifies considerably:

\[\frac{dN_j}{dt} = \int^V R_jdV = R_j \int^V dV = R_j V \]

In other words, as long as the reaction rate is uniform throughout the
volume of the reactor, we can say that:

\[G_j = R_j V\]

That gives a really useful result--the material balance for a well-mixed
batch reactor. We will use this balance extensively in reactor design
and analysis!

\[\frac{dN_j}{dt} = R_j V \]

    \hypertarget{simplification-for-a-constant-volume-process}{%
\subsubsection{Simplification for a Constant Volume
Process}\label{simplification-for-a-constant-volume-process}}

We cannot always guarantee that fluid density is constant in a batch
reactor, so we generally use the above equation since it allows us to
account for changes in the volume of reaction media that occur as the
reaction progresses. In some cases--gas phase reactions where there is
no change in moles or liquid phase reactions where there is a large
excess of solvent--we can assume the volume of the batch reactor is
constant. We will work through a useful simplification that applies only
for constant volume batch reactors. In addition to being widely used for
analysis of laboratory and kinetic data, this is a useful example
because it introduces the idea that you can manipulate differential
equations (and transform the differential variables therein) to suit the
system that you are considering. If we start with the general balance on
a well-mixed batch reactor:

\[\frac{dN_j}{dt} = R_j V \]

Without making any assumptions beyond perfect mixing, we can rearrange
as follows:

\[\frac{1}{V}\frac{dN_j}{dt} = R_j \]

Now, we'll assume that volume is time-independent (it does not change
with time). If reactor volume is constant with respect to time, we can
move \(1/V\) inside of the derivative:

\[\frac{d\frac{N_j}{V}}{dt} = R_j \]

We recognize that \(N_j/V\) is our standard definition for the molar
concentration of species \(j\), which gives:

\[\frac{dC_j}{dt} = R_j \]

This equation is frequently applied in the analysis of batch reactors.
Often, this is done without considering whether the constant volume
approximation is appropriate. One reason people do this is very likely
that \(\frac{dC_A}{dt} = R_j\) is an easy equation to remember, and
you've seen it elsewhere in, e.g., most of your chemistry courses. It is
very important that we remember that this result is only valid for a
constant-volume batch reactor. If our system experiences any sort of
change in density upon reaction, then we cannot use this simplified
result.

    \hypertarget{continuous-stirred-tank-reactor-cstr}{%
\subsection{Continuous Stirred Tank Reactor
(CSTR)}\label{continuous-stirred-tank-reactor-cstr}}

Next, we consider a CSTR. This is essentially a batch reactor where we
relax the constaints on material exchange across the boundaries to
create an open system. Specifically, a CSTR basically behaves like a
stirred batch wherein we allow flow of material into and out of the
reactor (See figure below). CSTRs are used frequently for liquid-phase
reactions, especially for those where mixing is beneficial to either
activity or selectivity (mixing is not always desirable).

\includegraphics{CSTR.png}

To derive a material balance on species \(j\) in a CSTR, we start with
the general balance equation:

\[\frac{dN_j}{dt} = F_{j,f} - F_j + G_j\]

Upon inspection (and consideration of the CSTR figure and operating
paradigm), we should conclude that none of the terms in the general
balance equation are guaranteed to be equal to zero. There is flow into
the system and flow out of the system, so the flow terms are non-zero.
There is also usually a reaction occuring, so we expect some
generation/consumption of species \(j\), which means that
\(G_j \neq 0\). Finally, unless we know specifically that the reactor is
at steady state, the quantity of each species may change with time, so
the accumulation term is non-zero. This idea tends to confuse students
at first because they are used to thinking about flow reactors at steady
state in material balence courses, but imagine a bucket with a hole in
it. At time = 0, the bucket is empty, and you turn on a faucet to fill
the bucket. The level of water will increase (i.e., there is
accumulation of water) until the system reaches a point where inflow and
outflow are balanced (there is no reaction when you're filling a bucket
with water, so \(G_j = 0\)). At that point, the amount of water in the
bucket stops changing, flow in equals flow out, there is no more net
accumulation of water, and the system is at steady state. But it takes
some time for this to happen! We will encouter a similar scenario any
time we start up or shut down a flow reactor--the flow reactor will
operate away from steady state during these times. For this notebook, we
will consider a homogeneous reaction occuring throughout the volume of
fluid, so we normalize reaction rates per unit volume. In this case, we
find that the relevant species balance for a CSTR is:

\[\frac{dN_j}{dt} = F_{j,f} - F_j + \int^V R_j dV\]

    \hypertarget{simplifying-the-cstr-balance-by-assuming-perfect-mixing}{%
\subsubsection{Simplifying the CSTR balance by assuming perfect
mixing\ldots{}}\label{simplifying-the-cstr-balance-by-assuming-perfect-mixing}}

Next, we assume that the reactor contents are perfectly mixed. Just like
in the batch reactor example, this means that temperature, pressure, and
composition inside of the reaction mixture are spatially uniform. They
do not change at all with position in the reactor, and if we take a
sample at one position, it will be identical to samples taken at all
other positions. One implication of the perfect mixing assumption is
that the reaction rate must be uniform throughout the reactor volume. In
this case, just as in a perfectly mixed batch reactor, we find:

\[\int^V R_j dV = R_j V\]

So our simplified balance on a perfectly mixed CSTR is:

\[\frac{dN_j}{dt} = F_{j,f} - F_j + R_j V\]

The well-mixed CSTR therefore represents the theoretical limit of a
\textbf{\emph{perfectly mixed flow reactor}}. In this course, we will
always assume that CSTRs are perfectly mixed.

    \hypertarget{simplifying-the-cstr-balance-by-assuming-steady-state}{%
\subsubsection{Simplifying the CSTR balance by assuming Steady
State\ldots{}}\label{simplifying-the-cstr-balance-by-assuming-steady-state}}

With flow reactors, we are frequently interested in their performance at
steady state. This happens after they are left on stream for a
sufficient period of time for dynamic (time dependent) processes to
decay to zero. At that point, the system reaches steady state. At steady
state, the system is time-independent. There are changes in temperature,
pressure, and/or species composition that occur between the inlet of the
reactor and the outlet of the reactor, but those changes will be the
same if we look at a steady state system now, 10 minutes in the future,
or 10 years in the future. The key is that once a system reaches steady
state, it becomes \textbf{\emph{time-independent}}. Once the system
becomes time-independent, there can be no accumulation of species. So,
when a CSTR reaches steady state, we find (for a perfectly mixed
reactor):

\[0 = F_{j,f} - F_j + R_j V\]

Usually, when we consider the design of CSTRs, we will be interested in
steady state performance, and we'll frequently use this equation.

    \hypertarget{other-implications-of-the-perfect-mixing-assumption-in-a-cstr}{%
\subsubsection{Other implications of the perfect mixing assumption in a
CSTR}\label{other-implications-of-the-perfect-mixing-assumption-in-a-cstr}}

Above, we pointed out that a perfectly mixed CSTR volume implies that
the temperature, pressure, and composition are uniform throughout the
reactor volume. This leads to some interesting and non-intuitive
behavior. We would expect that any flow reactor will accomplish a change
between the reactor inlet and the reactor outlet (otherwise, it has been
poorly designed). What this means is we should generally expect that the
temperature, pressure, and composition of the fluid entering a CSTR is
different from the temperature, pressure, and composition of the fluid
exiting a CSTR. We'll add this concept to our illustration of the CSTR:

\includegraphics{CSTRINOUT.png}

Generally, the above figure implies is that we should expect that, for
any CSTR, the volumetric flowrate (\(Q\)), the temperature (\(T\)), the
pressure (\(P\)), the species concentration (\(C_j\)), and/or the
species composition (\(x_j\), e.g., mole fractions) will be different in
the exit stream than they are in the feed stream. Here, we use a
subscript, \(f\), to denote ``feed'' conditions and no subscript to
denote exit conditions.

\hypertarget{consider-the-step-change}{%
\paragraph{Consider the step change}\label{consider-the-step-change}}

We now consider a subtle extension of the idea that temperature,
pressure, and composition are spatially uniform inside of a well-mixed
CSTR. One might imagine how the exit pipe from a CSTR is simply sampling
a part of the CSTR volume and removing it from the reactor. We have
already discussed how the assumption of perfect mixing implies that
everything inside of the CSTR is present at the same (\(T, P, x_j\)).
Since the exit pipe is basically removing a sample of the fluid from the
CSTR, so this means that the temperature, pressure, and composition in
the exit stream is exactly the same as the temperature, pressure, and
composition inside of the fluid phase within the reactor. We also know
that temperature, pressure, and composition in the fluid phase within
the reactor are the same everywhere inside of the reactor. Perfect
mixing then also implies that there is an entirely non-intuitive step
change between inlet conditions and exit conditions. This is illustrated
in the sketch below. This is how we should always think of a perfectly
mixed CSTR.

\includegraphics{CSTRStep.png}

    \hypertarget{tubular-reactor-plug-flow-reactor}{%
\subsection{Tubular Reactor (Plug Flow
Reactor)}\label{tubular-reactor-plug-flow-reactor}}

The final type of reactor that we'll consider (for now) is the Plug Flow
Reactor (PFR). These reactors are widely used for gas- or liquid phase
reactions in industry. A PFR is basically an empty tube through which
the reacting fluid is flowing---there is absolutely no mixing in a Plug
Flow Reactor. The core assumption for the PFR is that the fluid flows
through at a relatively high Reynolds number (\(\mathrm{Re} > 10^3\)).
This generally ensures that the fluid moves through the reactor in a
turbulent flow regime. This is convenient from a reactor modelling
perspective because it means that (if we consider an axial cross section
of the reactor), the velocity profile is flat, i.e., all fluid in that
cross section has an identical average velocity. Constrast this with a
Laminar Flow Regime where we assume there is a parabolic velocity
profile (See Figure Below).

\includegraphics{PFRvsLaminar.png}

The consequence of these assumptions is that everything moves through
through the PFR with the exact same velocity profile, so every molecule
we put in at the entrance of a PFR spends exactly the same amount of
time in the reactor. The way we think about the PFR is this: if we look
at a cross section of the tube, that entire cross section has spent an
identical amount of time in the reactor. So: everything \textbf{\emph{in
a single reactor cross section must be completely uniform in terms of
temperature, pressure, and composition}}. That said, as we move from one
cross section of the reactor to the next cross section travelling down
the length of the reactor, the temperature, pressure, and composition
will probably change. There is no mixing between cross sections, a PFR
behaves like an infinite series of isolated, well-mixed axial slices,
and we find that temperature, pressure, and composition vary as one
moves from one axial slice to the next down the length (through the
volume) of a PFR.

    \hypertarget{a-model-for-the-pfr}{%
\subsubsection{A model for the PFR}\label{a-model-for-the-pfr}}

Translating these ideas into a representative balance equation, as
discussed above, we will consider the PFR as being comprised of many
thin axial slices, all with an identical volume of \(\Delta V\).

\includegraphics{PFRDV.png}

Conceptually we will make each cross section very thin, so we are then
able to assume that there is no change in temperature, pressure, or
composition spatially throughout that entire cross section. Essentially,
we treat it like a very small CSTR! It is only when we move to the next
cross section that we would experience a change in temperature,
pressure, and composition. As we move axially down the length of a PFR,
we then see changes in temperature, pressure, and composition, but
\textbf{\emph{if we consider a single, isolated thin cross section}},
there is no variation in temperature, pressure, or composition with
respect to the radial or theta coordinate (i.e., everything in the cross
section is spatially uniform). With this concept in mind, we can propose
a mathematical description of this reactor.

    \hypertarget{focus-on-a-single-cross-section}{%
\subsubsection{Focus on a single cross
section\ldots{}}\label{focus-on-a-single-cross-section}}

We develop a PFR balance by considering a single cross sectional element
with a volume of \(\Delta V\):

\includegraphics{PFRsingleDV.png}

Inside of that cross section, everything is spatially uniform (by virtue
of having spent the same amount of time in the reactor), so we know that
this isolated cross section has a spatially uniform temperature,
pressure, and composition. It must also then have spatially uniform
reaction rates, \(r_i\), and spatially uniform species production rates,
\(R_j\). So, within that well-mixed cross section:

\[G_j = R_j\Delta V\]

For this derivation, we will assume that the PFR is operating at steady
state. This is not necessary, but we will only use the steady state
result in this class, so we will go ahead and apply the assumption from
the start since it makes the rest of the derivation more
straightforward. With these things in mind, we can write a balance on
that well-mixed cross section at steady state:

\[ 0 = F_j\big|_V - F_j\big|_{V+\Delta V} + R_j \Delta V\]

Next, we divide the whole equation by \(\Delta V\):

\[ 0 = \frac{F_j\big|_V - F_j\big|_{V+\Delta V}}{\Delta V} + R_j\]

And we now assume we have an infinite number of cross sections and let
the volume of each cross section approach zero, i.e.,
\(\Delta V \longrightarrow 0\). You may recognize this as the definition
of a derivative from your first semester calculus course. This gives:

\[ 0 = -\frac{dF_j}{dV} + R_j\]

In PFR analysis, we typically rearrange this equation to get:

\[\frac{dF_j}{dV} = R_j\]

Notice, unlike the general balance equation or material balances
developed for batch reactors and CSTRS, the PFR balance is actually an
\textbf{\emph{intensive}} balance because it is normalized per unit
volume of reactor. Inspection reveals that the dimensions on the left
and right hand side in a PFR balance equation are in moles per volume
per time. This equation gives us good intuition of how we think about a
PFR. It is a reactor that (generally) we will consider at steady state,
and we have to think about the way that composition (and temperature and
pressure) vary as a packet of fluid moves through the reactor from inlet
to exit. This is conceptually different from a batch reactor where there
is no variation with respect to position in the reactor, and we think
about changes that occur as a function of clock time.

Interestingly enough, the mathematics in PFR and Batch reactor problems
work out almost identically!


    % Add a bibliography block to the postdoc
    
    
    
\end{document}
